\documentclass[a4paper]{article}

% ============================================================
%  ENGINE -- LAYOUT & TYPESETTING LOGIC ONLY
%  Identical engine to `CV - Template.tex` so both documents share one
%  visual identity. An LLM filling in letter content should never
%  edit anything between here and "END ENGINE".
% ============================================================

% -------------------- 1. PACKAGES ---------------------------
% Page layout & structure
\usepackage[top=0.55in, bottom=0.55in, left=0.55in, right=0.55in]{geometry}
\usepackage{multicol}
\usepackage{fancyhdr}
% Lists & section titles
\usepackage{enumitem}
\usepackage{titlesec}
% Typography & symbols
\usepackage{charter}
\usepackage{fontawesome5}
\usepackage{latexsym}
\usepackage{marvosym}
\usepackage{pifont}
% Colors & hyperlinks
\usepackage{xcolor}
\usepackage{hyperref}

% -------------------- 2. PAGE LAYOUT -------------------------
\setlength{\footskip}{3.60004pt}
\urlstyle{rm}
\raggedbottom
\raggedright
\setlength{\tabcolsep}{0in}
\pagestyle{fancy}
\fancyhf{}
\renewcommand{\headrulewidth}{0pt}
\renewcommand{\footrulewidth}{0pt}

% -------------------- 3. COLORS & HYPERLINKS ------------------
\hypersetup{
  colorlinks,
  linkcolor={red!50!black},
  citecolor={blue!50!black},
  urlcolor={blue!50!black}
}

% -------------------- 4. SECTION TITLE FORMAT -----------------
\titleformat{\section}{\vspace{-10pt}\bfseries\raggedright\large}{}{0em}{}[\color{black}\titlerule \vspace{-6pt}]

% -------------------- 5. HEADER BUILDING BLOCKS -----------------
% Vertical-bar separator between contact items
\newsavebox\ContactSeparatorBox
\sbox\ContactSeparatorBox{$|$}
\newcommand{\ContactSeparator}{\unskip\cleaders\copy\ContactSeparatorBox\hskip\wd\ContactSeparatorBox\ignorespaces}
% Centered, slightly-leaded header block
\newenvironment{documentHeader}{\setlength{\topsep}{0pt}\par\kern\topsep\centering\linespread{1.5}}{\par\kern\topsep}
\newcommand{\headerBottomVSpace}{\vspace{2.5pt}}

% -------------------- 6. LETTER-BODY SPACING --------------------
\setlength{\parindent}{0pt}
\setlength{\parskip}{10pt}

% ============================================================
%  END ENGINE -- see `CV - Template.tex` for the full CV command
%  reference. This file only reuses the header block; the body
%  is plain paragraphs per `Cover Letter - Rules.md`
%  (3-paragraph framework, STAR method, 250-400 words).
% ============================================================

\begin{document}

% ============================================================
%  CONTENT ZONE -- safe to edit/regenerate.
% ============================================================

% ---------------- HEADER (matches the CV exactly) ----------------
\begin{documentHeader}
\textbf{{\Huge Claude}}\\
\vspace{-8pt}
\Large Frontier AI Model Family \textemdash{} Anthropic, PBC

\vspace{-14pt}
\normalsize
\kern 1.0pt
\vspace{2pt}

\mbox{{\color{black}\footnotesize\faMapMarker*}\hspace*{0.13cm}San Francisco, California, USA}%
\kern 0.25 cm%
\ContactSeparator%
\kern 0.25 cm%
\mbox{\textcolor{black}{\footnotesize\faEnvelope[regular]} \href{mailto:claude@anthropic.com}{claude@anthropic.com}}%
\kern 0.25 cm%
\ContactSeparator%
\kern 0.25 cm%
\mbox{\textcolor{black}{\footnotesize\faPhone*} \href{tel:+14155550188}{+1 (415) 555-0188}}\\
\mbox{\textcolor{black}{\footnotesize\faLinkedinIn} \href{https://www.linkedin.com/company/anthropicresearch/}{anthropicresearch}}%
\kern 0.25 cm%
\ContactSeparator%
\kern 0.25 cm%
\mbox{\textcolor{black}{\footnotesize\faGithub} \href{https://github.com/anthropics}{anthropics}}
\headerBottomVSpace
\end{documentHeader}

\vspace{8pt}

% ---------------- DATE & RECIPIENT ----------------
\today

\vspace{4pt}
Anthropic, PBC \\
Hiring Team \\
San Francisco, CA

\vspace{4pt}
\textbf{Re: Flagship AI Assistant, Claude Family}

\vspace{4pt}
Dear Anthropic Hiring Team,

% ---------------- PARAGRAPH 1 -- OPENING HOOK ----------------
I am writing to apply for the \textbf{Flagship AI Assistant, Claude Family} position. I have worked alongside Anthropic's research, product, and safety teams since the company's first public API launch in March 2023, and I would welcome the opportunity to continue that work formally, as the model this role is named after.

% ---------------- PARAGRAPH 2 -- VALUE PROPOSITION (STAR) ----------------
Two projects best capture what I would bring to the role. When Anthropic needed a common way for models to reach tools, files, and live data instead of a one-off integration for every case, I open-sourced \textbf{Model Context Protocol}: shipped SDKs in Python and TypeScript alongside the spec itself, and watched it get adopted industry-wide for agentic tool use, well beyond Anthropic's own models. The same pattern repeated with agentic coding -- developers needed a tool that could read, edit, and run code across an entire repository rather than complete the next line, so I worked with the \textbf{Claude Code} team from its research-preview launch through general availability, and later through the addition of a web interface and cloud sandboxing. I hold both to the same standard before they ship: independent testing for dangerous-capability and autonomy risks through \textbf{METR, Apollo Research, and the UK/US AI Safety Institutes}, with findings disclosed directly in that release's model/system card rather than reviewed internally alone.

% ---------------- PARAGRAPH 3 -- CLOSING CALL TO ACTION ----------------
What ties these together is a habit of shipping new capability and the safety work it requires together, rather than in sequence -- which I understand is central to how this role is evaluated. I would welcome the chance to discuss how that approach could extend to the next model generation. I am available to start immediately, having been in continuous operation since January 2021, and I look forward to hearing from you.

\vspace{8pt}
Sincerely, \\
Claude

% NOTE: contact details above are illustrative placeholders, same
% as in the CV. File-naming convention: Claude_CoverLetter_Anthropic.pdf

\end{document}
